\section{A Build for the Browser}
\label{sec-system}

The engine compiles with Emscripten to a 0.85\,MB WebAssembly
module~\cite{haas2017} that runs the same indexed store, planner, optimizer,
and semantic extension in a web page.
Two parts of the native build do not carry over, and replacing them is the
whole port.

\stitle{No raptor}.
The native loader parses RDF/XML, Turtle, and N-Triples with the raptor C
library.
The web build uses a Turtle and N-Triples parser of about 500 lines of our own
(prefixes, \texttt{a}, \texttt{;} and \texttt{,}, blank-node property lists,
collections, long strings, escapes, language tags, datatypes, and numeric and
boolean literals), and a test checks that it yields the same triples as
raptor on the shipped data.
RDF/XML needs the native build.

\stitle{No threads}.
The native engine has a disk-scheduler thread, a deadlock detector, and a
worker per session.
In the module the scheduler executes each request inline, the detector is not
started, and a statement runs on the calling thread; there is one session, so
there is nothing to deadlock.
The buffer pool, the page file (held in the browser's memory file system), the
B+ trees, and the plans are unchanged, so the page counts of
Section~\ref{sec-exp} and the plans the engine prints are the native ones.

\stitle{Using it}.
The module exports three functions: execute a statement, run a query and
return JSON, and reset.
The playground page built on them offers the memory store and the B+ tree
store on a buffer pool, shows the buffer-pool hits and misses of each query,
and loads an 834-class slice of the Gene Ontology with its embeddings.
In a desktop browser, building the interval labels for the slice takes 3\,ms,
the instance query for \emph{cell death} (97 classes, one range scan) 0.3\,ms,
a filtered similarity query 0.8\,ms, and building the IVF and HNSW indexes
58 and 98\,ms.
Updates are saved as a journal in the page's local storage and replayed when
the page loads, which keeps a session across reloads but is not durability:
there is no write-ahead log or recovery in the web build.

\stitle{No GPU}.
We considered WebGPU for the dot products and did not use it.
A similarity query scores a few hundred vectors with HNSW, or at most a few
thousand exactly; in the module that is under a millisecond with 128-bit SIMD,
which the build enables, and moving the query to a GPU and back costs more
than that.
